\corrigesChapitre{10}{Interactions et forces}
\par\noindent\begin{minipage}{\linewidth}
\begin{corrige}{1}{Choisir le système}
\begin{questions}
\item Le ballon, si c’est l’objet annoncé.
\item Oui ; son action sur le ballon doit être examinée.
\end{questions}
\end{corrige}
\end{minipage}
\par\noindent\begin{minipage}{\linewidth}
\begin{corrige}{2}{Contact ou distance}
\begin{questions}
\item Table : contact ; Terre : distance.
\item Non : l’interaction magnétique peut s’exercer à distance.
\end{questions}
\end{corrige}
\end{minipage}
\par\noindent\begin{minipage}{\linewidth}
\begin{corrige}{3}{Quatre caractéristiques}
\begin{questions}
\item Point d’application, direction, sens et valeur.
\item Le newton, symbole N.
\end{questions}
\end{corrige}
\end{minipage}
\par\noindent\begin{minipage}{\linewidth}
\begin{corrige}{4}{L’appareil adapté}
\begin{questions}
\item Un dynamomètre.
\item Pour éviter de le détériorer et utiliser une plage de mesure adaptée à la valeur attendue.
\end{questions}
\end{corrige}
\end{minipage}
\par\noindent\begin{minipage}{\linewidth}
\begin{corrige}{5}{Deux sens}
\begin{questions}
\item Oui : leurs droites supports sont verticales.
\item Non : les sens sont opposés.
\end{questions}
\pcFigureEx{figures/cor-direction-sens.pdf}{Même direction verticale ; sens opposés. Les flèches représentent ici deux forces de même valeur, à titre d’exemple.}
\end{corrige}
\end{minipage}
\par\noindent\begin{minipage}{\linewidth}
\begin{corrige}{6}{Dessiner à l’échelle}
\begin{questions}
\item $6/2=3\unite{cm}$.
\item $1/2=0{,}5\unite{cm}$.
\end{questions}
\par\smallskip\noindent\textbf{Pour comprendre.}\enspace Chaque segment de 1 cm représente 2 N. La longueur vaut donc $\ell=F/2$ en cm : 3 cm pour 6 N et 0,5 cm pour 1 N. Mesurer du point de départ à la pointe, sans changer l’échelle entre les deux flèches.\par
\pcFigureExact{figures/cor-forces-echelle.pdf}{Longueurs à mesurer du point d’application à la pointe. Échelle conservée pour une impression à 100 \%.}
\end{corrige}
\end{minipage}
\par\noindent\begin{minipage}{\linewidth}
\begin{corrige}{7}{Lire une longueur}
\begin{questions}
\item $4\times0{,}5=2\unite N$.
\item Non : cela change seulement la longueur choisie pour la représentation.
\end{questions}
\par\smallskip\noindent\textbf{Pour comprendre.}\enspace Chaque centimètre vaut 0,5 N, donc $F=4\times0{,}5=2\unite N$. Une autre échelle change le dessin, pas la force physique représentée.\par
\pcFigureExact{figures/cor-force-4cm.pdf}{Quatre segments de 1 cm représentent chacun 0,5 N.}
\end{corrige}
\end{minipage}
\par\noindent\begin{minipage}{\linewidth}
\begin{corrige}{8}{Un objet immobile}
\begin{questions}
\item Non : il peut subir son poids.
\item D’autres forces, par exemple celle d’un support, peuvent compenser cette action dans le modèle étudié.
\end{questions}
\end{corrige}
\end{minipage}
\par\noindent\begin{minipage}{\linewidth}
\begin{corrige}{9}{Le livre sur la table}
\begin{questions}
\item Le livre est relié à la Terre par une interaction à distance et à la table par une interaction de contact.
\item L’action de la Terre vers le bas et celle de la table vers le haut, dans cette situation simple.
\end{questions}
\par\smallskip\noindent\textbf{Pour comprendre.}\enspace Choisir d’abord le livre comme système. Le diagramme recense Terre et table. Le bilan des forces garde seulement les actions \emph{sur le livre}, et non celles exercées par le livre sur ses partenaires.\par
\pcFigureEx{figures/cor-livre-forces.pdf}{Le diagramme recense les partenaires ; les flèches représentent les actions sur le livre, dans le modèle immobile à deux forces.}
\end{corrige}
\end{minipage}
\par\noindent\begin{minipage}{\linewidth}
\begin{corrige}{10}{Une balle qui tourne}
\begin{questions}
\item Oui : la vitesse peut changer de direction à valeur constante.
\item Un mouvement circulaire uniforme, par exemple une balle retenue dans un mouvement circulaire adapté ; il ne faut pas réaliser de lancer dangereux.
\end{questions}
\end{corrige}
\end{minipage}
\par\noindent\begin{minipage}{\linewidth}
\begin{corrige}{11}{Deux objets, deux actions}
\begin{questions}
\item Oui : il s’agit d’une interaction réciproque.
\item Non : une force s’applique à la boîte, l’autre à la main. Le bilan de la boîte n’inclut que les forces qui s’appliquent à la boîte.
\end{questions}
\par\smallskip\noindent\textbf{Pour comprendre.}\enspace La main agit sur la boîte et la boîte agit sur la main. Les receveurs diffèrent : le bilan mécanique de la boîte contient l’action de la main, mais pas l’action reçue par la main.\par
\pcFigureEx{figures/cor-main-boite.pdf}{Les deux actions de l’interaction s’appliquent à deux systèmes différents : elles ne se compensent pas sur la seule boîte.}
\end{corrige}
\end{minipage}
\par\noindent\begin{minipage}{\linewidth}
\begin{corrige}{12}{Un bilan incomplet}
\begin{questions}
\item L’action de la Terre, le poids.
\item Il pourrait croire qu’une seule force vers le haut suffit à expliquer l’immobilité, alors que les actions doivent être comparées.
\end{questions}
\par\smallskip\noindent\textbf{Pour comprendre.}\enspace Le poids n’est pas supprimé par l’immobilité. Dans le modèle à deux forces, la tension du fil compense le poids : même direction et même valeur, sens opposés.\par
\pcFigureEx{figures/cor-suspension.pdf}{Il faut ajouter le poids ; dans le modèle immobile à deux forces, les valeurs sont égales et les sens opposés.}
\end{corrige}
\end{minipage}
\par\noindent\begin{minipage}{\linewidth}
\begin{corrige}{13}{Une lecture discutable}
\begin{questions}
\item Soustraire le décalage : 2,0 N.
\item Régler le zéro correctement avant la mesure et vérifier la stabilité, plutôt que corriger sans contrôle.
\end{questions}
\end{corrige}
\end{minipage}
\par\noindent\begin{minipage}{\linewidth}
\begin{corrige}{14}{Une flèche sans légende}
\begin{questions}
\item Aucun nécessairement : leurs échelles peuvent être différentes.
\item Comparer échelle, valeur, point d’application, direction, sens et système étudié.
\end{questions}
\end{corrige}
\end{minipage}
\par\noindent\begin{minipage}{\linewidth}
\begin{corrige}{15}{Changer le système}
\begin{questions}
\item Elle devient une interaction interne à l’ensemble objet-fil.
\item Il faut notamment considérer l’action du support sur le fil et celle de la Terre sur l’ensemble.
\end{questions}
\pcFigureEx{figures/cor-systemes.pdf}{Une action devient interne quand les deux partenaires sont inclus dans le système choisi.}
\end{corrige}
\end{minipage}
\par\noindent\begin{minipage}{\linewidth}
\phantomsection\label{cor-defi-10}\pcSolutionDefi{10}{Une suspension à expliquer}\begin{questions}
\item Le système est l’objet ; il interagit avec le fil par contact et avec la Terre à distance.
\item Deux flèches verticales opposées de $3/1{,}5=2\unite{cm}$, si le modèle à deux forces est valable.
\item L’immobilité n’implique pas l’absence de forces. Ici, leurs effets se compensent ; leurs valeurs ne sont pas nulles.
\end{questions}
\pcFigureExact{figures/cor-suspension-3N.pdf}{Chaque flèche mesure 2 cm à l’impression à 100 \%. Les deux forces s’appliquent au même objet et se compensent dans ce modèle.}
\end{minipage}
\par\noindent\begin{minipage}{\linewidth}
\phantomsection\label{cor-activite-10}\pcSousTitre{Pistes pour l’investigation}
Le système étudié est la masse suspendue. Elle interagit avec la Terre à distance et avec le dynamomètre par contact. À l’équilibre dans un modèle à deux forces, l’action de suspension et le poids ont la même direction, des sens opposés et des valeurs égales. L’instrument mesure l’action selon son axe ; son alignement et son zéro doivent être contrôlés.\par\end{minipage}
